\chapter{Limitações e conclusões}
\section{Alcance da avaliação}
A investigação concentrou-se em uma área de distribuição e em um horizonte de 117 dias civis (01/08 a 25/11/2026), com avaliação fora do período de ajuste em dois cortes distintos (19/11 e 25/11/2026). Há dependência entre janelas da mesma UC e heterogeneidade entre circuitos. UCs sem alimentador e instalações com telemetria insuficiente seguem percursos distintos de análise. A cobertura incompleta do vínculo a transformadores exigiu recorrer à grade geográfica de aproximadamente 250~m como nível intermediário de corroboração espacial. O sistema também depende da disponibilidade de alarmes e da cadência efetiva de cada canal.

A avaliação experimental combinou inferência OOT sobre dados reais, auditorias de sensibilidade com modelo congelado, validação temporal multi-corte e perturbações fisicamente coerentes sobre cópias isoladas de séries reais. Esse protocolo demonstra estabilidade estatística, sensibilidade controlada e complementaridade entre especialistas, mas não substitui rótulos externos de inspeção: a confirmação pericial de fraude, o rendimento efetivo de campo (\emph{yield}), a energia recuperada e o retorno econômico realizado permanecem fora do escopo experimental por ausência de retorno de OS e TOI. Os pesos do comitê, os limiares do roteador hierárquico e as ponderações de $S_{tamper}$ permanecem escolhas heurísticas de engenharia validadas por sensibilidade, cuja calibração ótima requer desfechos independentes em piloto prospectivo.

A auditoria da execução de referência (\texttt{v1}, 25/11) identificou três ajustes implementados na variante \texttt{v2} sem sobrescrever os artefatos originais. Primeiro, o filtro textual do bônus $+0{,}15$ alcançava o desequilíbrio de corrente em 46,73\% da coorte, sendo restringido na \texttt{v2} às condições de tensão e reversão verificadas no ensaio de sensibilidade do bônus (0,19\% da coorte). Segundo, a leitura direta das partições de \texttt{alarm\_events} substituiu a busca por arquivo único ausente, ativando $S_{\text{tamper}} \ge 0{,}70$ em 1.377 UCs no corte de 19/11. Terceiro, o requisito adicional de 18 dias confiáveis bloqueou os silêncios consecutivos de 15 dias que contornavam a média dos dias presentes no teste controlado de lacunas.

\section{Conclusões e continuidade}
A reconstrução da vigência \texttt{UC--NIO}, aliada ao saneamento físico das grandezas elétricas (corrente sob carga $\ge 1{,}0\ \mathrm{A}$, detecção automática de fases ativas e cobertura por cadência nativa), alterou substancialmente a interpretabilidade dos casos. Na UC \texttt{908***00}, reconhecer a continuidade do equipamento relacionou a queda de FA ao histórico anterior e à injeção reversa mantida até 25/11/2026. Na UC \texttt{291***96}, a progressão do registrador reverso após a retirada declarada e o silêncio prolongado de telemetria motivaram a auditoria de vigência cadastral e a retenção segura em quarentena. Na UC \texttt{685***20}, a reconciliação entre grade legada e cadência nativa fechou a divergência entre 50,00\% e 99,49\%, mantendo o caso no topo prioritário até 25/11.

No plano algorítmico, os experimentos consolidados sustentam quatro conclusões técnicas objetivas:
\begin{enumerate}
\item \textbf{Estabilidade e integridade da escala:} a normalização quantílica congelada ($P_{01}$--$P_{99}$) garantiu deslocamento exatamente nulo ($0{,}000000$) sob perturbação de 1\% de outliers, frente a $0{,}2291$ na escala dinâmica por lote. A auditoria das correlações confirmou alinhamento positivo entre todos os seis detectores (\texttt{IF} $\times$ \texttt{ECOD} $= +0{,}9626$ nos elegíveis), com diversidade efetiva aportada por LOF (180--187 substituições no Top-2.000) e PCA residual (166 substituições), além de robustez a variações de $\pm 20\%$ nos pesos (Jaccard Top-2.000 entre $0{,}9724$ e $0{,}9940$).
\item \textbf{Efeito marginal das visões na coorte real e nas perturbações:} na coorte real de 19/11 (Base A, $K = 2.000$), o filtro de cobertura eliminou a entrada de unidades em quarentena no Top-2.000 (de 138 em V0 e 257 em V1 para 0), elevou de 134 para 198 os casos com queda sustentada no CUSUM e produziu 282 trocas com o roteamento MoE por alimentador. No conjunto com perturbações conhecidas sobre 100 séries reais (Base B, $N = 800$, $K = 400$), a combinação das visões elevou o F1 em perturbações elétricas de 68,25\% (V0) para 76,00\% (V5), localizou quedas de $-50\%$ com erro mediano de 0,0~dia, detectou 52,0\% das rotações horárias de $+4\ \mathrm{h}$ (frente a 32,0\% no escore diário) e reclassificou 16/16 quedas simultâneas de circuito como evento de rede.
\item \textbf{Retenção entre cortes e ajustes da variante \texttt{v2}:} entre 19/11 e 25/11/2026, permaneceram no Top-2.000 89,8\% das unidades da referência \texttt{v1} (1.797/2.000; $\rho = 0{,}9688$) e 88,0\% da variante \texttt{v2} (1.760/2.000; $\rho = 0{,}9484$), enquanto a carteira completa \texttt{v2} reteve 91,5\% das candidatas iniciais (3.933/4.298). A restrição do bônus $+0{,}15$ às regras de tensão e reversão reduziu seu acionamento de 46,73\% para 0,19\% da coorte mantendo 92,4\% do Top-2.000, o critério mínimo de 18 dias confiáveis bloqueou 100\% dos silêncios consecutivos de 15 dias na Base B (vazamento elétrico de 31,0\% para 0,0\%) e a leitura particionada de alarmes identificou 1.377 UCs em 19/11 e 1.375 em 25/11 com $S_{\text{tamper}} \ge 0{,}70$.
\item \textbf{Engenharia voltada à operação da distribuidora:} o firewall programático de subespaço elétrico, o contraste espacial em três camadas para RF Mesh (Posto $\to$ grade 250~m $\to$ Alimentador), o controle de sessões Oracle via \texttt{QueuePool} limitado e o versionamento SHA-256 tornam reconstituível todo o percurso entre a leitura de 5/10 minutos no medidor e a fila priorizada entregue ao analista.
\end{enumerate}

Como etapa seguinte, propõe-se executar o piloto prospectivo aleatorizado delineado no Capítulo~\ref{cap:operacionalizacao} para vincular as listas congeladas da variante \texttt{v2} aos desfechos independentes de OS e TOI. Somente com esses registros será possível ajustar pesos supervisionados e estimar, sob capacidade fixa de inspeção, a taxa de confirmação, a antecedência em relação à carteira convencional e o saldo econômico líquido.
